\section{Conjuntos abertos}\index{noções topológicas}
Nesta seção, discutiremos sucintamente a noção de conjunto aberto de um espaço métrico, com enfoque nos espaços $\mathbb R^n$.

Ao longo desta seção, $(M,d)$ denotará um espaço métrico arbitrário.

Aconselhamos o leitor a ter sempre em mente o exemplo de $\mathbb R^n$ com a distância euclidiana.

\begin{definition}\index{conjunto aberto}
    Seja $A$ um subconjunto de $M$.
    Dizemos que $A$ é \emph{aberto} se, para todo $p\in A$, existe $r>0$ tal que $B(p, r)\subseteq A$.
\end{definition}

A intuição de conjunto aberto é a seguinte: para que $A$ seja aberto, é necessário que cada um dos seus pontos possua uma ``folga'' dentro de $A$.

Observe que o raio $r$ da definição pode depender do ponto $p$.
Não se exige que exista um único raio positivo que funcione simultaneamente para todos os pontos de $A$.

\subimport{}{abertos-figuraAbertos}

Como um primeiro exemplo, verificaremos que bolas abertas são abertas, justificando esta nomenclatura.

\begin{proposition}
    Sejam $(M,d)$ um espaço métrico, $p \in M$ e $r>0$.
    Então, a bola aberta $B(p, r)$ é um conjunto aberto.
\end{proposition}
\begin{proof}
    Seja $q \in B(p, r)$ arbitrário.
    Note que $d(p, q) < r$.
    Intuitivamente, $r-d(p,q)$ mede a margem que ainda sobra na desigualdade $d(p,q)<r$ (ver Figura~\ref{fig:abertosBolaAberta}).

    Seja $s = r - d(p, q) > 0$.
    Veremos que $B(q, s) \subseteq B(p, r)$.

    De fato, seja $x \in B(q, s)$, ou seja, $d(q, x) < s$.
    Então:
    \begin{align*}
        d(p, x) &\leq d(p, q) + d(q, x) \\
        &< d(p, q) + s = d(p, q) + (r - d(p, q)) = r.
    \end{align*}

    Assim, $x \in B(p, r)$ e concluímos que $B(q, s) \subseteq B(p, r)$.

    Como $q \in B(p, r)$ foi arbitrário, concluímos que para todo $q \in B(p, r)$ existe $s>0$ tal que $B(q, s)\subseteq B(p, r)$.
    Logo, pela definição, $B(p, r)$ é aberto.
\end{proof}

\subimport{}{abertos-figuraBolaAberta}

Algumas propriedades básicas dos conjuntos abertos são as seguintes.

\begin{proposition}\index{conjunto aberto!propriedades}\label{prop:propriedadesAbertos}
    Seja $(M,d)$ um espaço métrico.
    Sobre os abertos de $M$, temos as seguintes propriedades.
    \begin{enumerate}[label=(\roman*)]
        \item $M$ e o conjunto vazio $\emptyset$ são abertos.
        \item A união de qualquer coleção de conjuntos abertos é aberta.\index{conjunto aberto!união}
        \item A interseção de dois conjuntos abertos é aberta.\index{conjunto aberto!interseção finita}
    \end{enumerate}
\end{proposition}
\begin{proof}
    (i) O conjunto $M$ é aberto porque, para todo $p \in M$, temos que $B(p, 1) \subseteq M$.
    O conjunto vazio $\emptyset$ é aberto por vacuidade: do contrário, haveria um ponto $p \in \emptyset$ tal que, para todo $r>0$, $B(p, r)\not\subseteq \emptyset$, o que é absurdo, pois não existem pontos no conjunto vazio.

    (ii) Seja $\{A_i\}_{i \in I}$ uma coleção indexada de conjuntos abertos de $M$.
    Seja $A = \bigcup_{i \in I} A_i$.
    Seja $p \in A$ arbitrário.
    Então, existe $j \in I$ tal que $p \in A_j$.
    Como $A_j$ é aberto, existe $r>0$ tal que $B(p, r) \subseteq A_j\subseteq A$, e, portanto, $B(p, r) \subseteq A$.
    Assim, $A$ é aberto.

    (iii) Sejam $A_1$ e $A_2$ conjuntos abertos de $M$.
    Seja $p \in A_1\cap A_2$ arbitrário.
    Existe $r_1>0$ tal que $B(p, r_1) \subseteq A_1$ e existe $r_2>0$ tal que $B(p, r_2) \subseteq A_2$.
    
    Seja $r = \min\{r_1, r_2\}$, que é positivo, pois $r_1$ e $r_2$ são positivos.
    Então, $B(p, r) \subseteq B(p, r_1) \subseteq A_1$ e $B(p, r) \subseteq B(p, r_2) \subseteq A_2$, ou seja, $B(p, r) \subseteq A_1\cap A_2$.
    Assim, $A_1\cap A_2$ é aberto.
\end{proof}
\subimport{}{abertos-figuraIntersecaoAbertos}

\begin{corollary}
    Sejam $(M,d)$ um espaço métrico e $A_1, \dots, A_n$ conjuntos abertos de $M$, com $n\geq 1$.
    Então a interseção $\bigcap_{i=1}^n A_i$ é um conjunto aberto de $M$.
\end{corollary}
\begin{proof}
    Procedemos por indução em $n$.

    O caso $n=1$ é trivial, pois a referida interseção é exatamente $A_1$, que é aberta por hipótese.

    Suponha que a hipótese é válida para algum $n\geq 1$.
    Veremos que a hipótese é válida para $n+1$.

    De fato, temos que
    \[
        \bigcap_{i=1}^{n+1} A_i = \left(\bigcap_{i=1}^n A_i\right)\cap A_{n+1}.
    \]
    Pela hipótese de indução, $\bigcap_{i=1}^n A_i$ é aberto.
    Além disso, $A_{n+1}$ é aberto por hipótese.
    Logo, $\bigcap_{i=1}^{n+1} A_i$ é a interseção de dois conjuntos abertos e, portanto, é aberto pelo item (iii) da Proposição~\ref{prop:propriedadesAbertos}.
\end{proof}

O corolário abaixo estabelece uma caracterização dos conjuntos abertos em um espaço métrico.

\begin{corollary}\label{corolario:abertos-uniao-bolas}
    Seja $M$ um espaço métrico e seja $A\subseteq M$.
    São equivalentes:
    \begin{enumerate}[label=(\roman*)]
        \item $A$ é aberto;
        \item Existe uma família $(r_p)_{p\in A}$ de números reais positivos tal que
        \[
            A = \bigcup_{p\in A} B(p, r_p).
        \]
        \item $A$ é a união de uma coleção de bolas abertas.
    \end{enumerate}
\end{corollary}
\begin{proof}
    $(i)\rightarrow (ii)$.
    Se $A$ é um conjunto aberto, então, para cada $p \in A$, existe $r_p>0$ tal que $B(p, r_p)\subseteq A$.
    Como $B(p,r_p)\subseteq A$ para todo $p\in A$, segue que $\bigcup_{p\in A} B(p, r_p) \subseteq A$.

    Temos ainda que $A \subseteq \bigcup_{p\in A} B(p, r_p)$, pois, dado $p \in A$, temos $p \in B(p, r_p)$.
    Logo, $A = \bigcup_{p\in A} B(p, r_p)$.

    $(ii)\rightarrow (iii)$ é imediato, pois se $A = \bigcup_{p\in A} B(p, r_p)$, então $A$ é a união de uma coleção de bolas abertas.

    $(iii)\rightarrow (i)$ decorre do item (ii) da Proposição~\ref{prop:propriedadesAbertos} e do fato de que bolas abertas são abertas.
\end{proof}
\begin{example}
    A interseção de infinitos conjuntos abertos não precisa ser aberta.
    Em $\mathbb R$, com a distância usual, temos
    \[
        \bigcap_{n=1}^{\infty} \left]-\frac{1}{n},\frac{1}{n}\right[ = \{0\}.
    \]
    Porém, $\{0\}$ não é aberto em $\mathbb R$, pois toda bola aberta centrada em $0$ contém pontos diferentes de $0$.
\end{example}

\begin{example}\label{exemplo:intervalos-abertos-sao-abertos}
    Em $\mathbb R$, com a distância usual, cada intervalo aberto $]a, b[$, com $a<b$, é um conjunto aberto.
    Para ver isso, podemos escrevê-lo como uma bola aberta.
    
    Seu ponto médio, $x=\frac{a+b}{2}$, é um candidato natural a ser o centro da bola aberta.

    Para acharmos um candidato a raio, primeiro vamos medir a distância entre $x$ e $a$.
    Temos que $d(x, a) = \left|\frac{a+b}{2} - a\right| = \frac{b-a}{2}$.
    Assim, seja $r = \frac{b-a}{2}$.

    Afirmo que $]a, b[=B(x, r)$.

    Com efeito, dado $c \in \mathbb R$, temos que:
    \begin{align*}
        c \in B(x, r) 
        &\Leftrightarrow |c-x| < r \\
        &\Leftrightarrow x-r<c<x+r \\
        &\Leftrightarrow \frac{a+b}{2}-\frac{b-a}{2}<c<\frac{a+b}{2}+\frac{b-a}{2} \\
        &\Leftrightarrow a<c<b \\
        &\Leftrightarrow c \in ]a, b[.
    \end{align*}

    Assim, $]a, b[=B(x, r)$ é um conjunto aberto.
\end{example}

\subimport{}{abertos-figuraIntervalosAbertos}

\begin{example}
    O intervalo aberto $]a, +\infty[\subseteq \mathbb R$, com $a \in \mathbb R$, é um subconjunto aberto de $\mathbb R$.
    
    Para ver isso, basta notar que $]a, +\infty[ = \bigcup_{b>a} ]a, b[$ é uma união de conjuntos abertos, logo é aberto.

    Analogamente, o intervalo aberto $]-\infty, a[$ é um subconjunto aberto de $\mathbb R$.
    De fato, $]-\infty, a[$ é uma união de conjuntos abertos, e, portanto, aberto, pois
    \[
        ]-\infty, a[ = \bigcup_{b<a} ]b, a[.
    \]
\end{example}

A noção de conjunto aberto depende do espaço ambiente.

\begin{example}
    Em $\mathbb R$, com a distância usual, o conjunto $[0,1[$ não é aberto.
    De fato, $0\in [0,1[$, mas, para todo $r>0$, a bola $B_{\mathbb R}(0,r)=]-r,r[$ contém pontos negativos, como $-r/2$.
    Portanto, nenhuma bola aberta de $\mathbb R$ centrada em $0$ está contida em $[0,1[$.

    Por outro lado, no subespaço $A=[0,+\infty[\subseteq\mathbb R$, com a métrica herdada, temos
    \[
        B_A(0,1)=B_{\mathbb R}(0,1)\cap A=[0,1[.
    \]
    Assim, $[0,1[$ é aberto no subespaço $A$, embora não seja aberto em $\mathbb R$.
\end{example}

Note também que se $M$ é um espaço métrico e $A\subseteq M$ é considerado como subespaço de $M$, então $A$ é aberto com relação ao próprio espaço métrico $A$, mesmo que não seja aberto em $M$.
De fato, basta aplicar o item (i) da Proposição~\ref{prop:propriedadesAbertos} ao espaço métrico $A$.

A relação entre os abertos de um subespaço e os do espaço ambiente é a seguinte.

\begin{proposition}
    Sejam $(M,d)$ um espaço métrico, $N\subseteq M$ um subespaço de $M$ e $A\subseteq N$.
    Então $A$ é aberto em $N$ se, e somente se, existe um conjunto aberto $B$ de $M$ tal que $A=B\cap N$.
\end{proposition}
\begin{proof}
    Suponha que $A$ é aberto em $N$.
    Do Corolário~\ref{corolario:abertos-uniao-bolas}, existe uma família $(r_p)_{p\in A}$ de números reais positivos tal que
    \[
        A = \bigcup_{p\in A} B_N(p, r_p).
    \]
    Assim,
    \[
        A = \bigcup_{p\in A} \bigl(B_M(p, r_p)\cap N\bigr) = \bigcup_{p\in A} \left(B_M(p, r_p)\cap N\right) = \left(\bigcup_{p\in A} B_M(p, r_p)\right)\cap N.
    \]
    Tomando $B = \bigcup_{p\in A} B_M(p, r_p)$, segue a tese, pois $B$ é aberto em $M$ e $A = B\cap N$.

    Reciprocamente, suponha que existe um conjunto aberto $B$ de $M$ tal que $A=B\cap N$.
    Seja $p\in A$.
    Então $p\in B$ e $p\in N$.
    Como $B$ é aberto em $M$, existe $r>0$ tal que
    \[
        B_M(p,r)\subseteq B.
    \]
    Logo,
    \[
        B_N(p,r)=B_M(p,r)\cap N\subseteq B\cap N=A.
    \]
    Portanto, $A$ é aberto em $N$.
\end{proof}

\begin{knowMore}
    A noção de espaço métrico pode ser generalizada para a noção de espaço topológico.

    Um \emph{espaço topológico} é, por definição, um par $(X, \tau)$, em que $X$ é um conjunto e $\tau$ é uma coleção de subconjuntos de $X$ tal que:
    \begin{enumerate}[label=(\roman*)]
        \item $X$ e $\emptyset$ pertencem a $\tau$;
        \item A união de qualquer coleção de conjuntos pertencentes a $\tau$ pertence a $\tau$;
        \item A interseção de dois conjuntos pertencentes a $\tau$ pertence a $\tau$.
    \end{enumerate}
    Nesse contexto, os conjuntos pertencentes a $\tau$ são chamados de \emph{conjuntos abertos} de $X$.
    Assim, pela Proposição~\ref{prop:propriedadesAbertos}, todo espaço métrico $(M,d)$ determina um espaço topológico $(M,\tau_d)$, em que $\tau_d$ é a coleção dos subconjuntos abertos de $M$.
\end{knowMore}

\subsection*{Exercícios}

\begin{exercise}
    Em $\mathbb R$, com a distância usual, determine quais dos seguintes subconjuntos são abertos.
    Justifique suas respostas.
    \begin{enumerate}[label=(\alph*)]
        \item $]-2,3[$;
        \item $[0,1[$;
        \item $\mathbb R\setminus\{0\}$;
        \item $\mathbb Z$;
        \item $\displaystyle \bigcup_{n=1}^{\infty} ]n,n+1[$.
    \end{enumerate}
\end{exercise}

\begin{exercise}
    Em $\mathbb R^n$, com a distância euclidiana, mostre que o conjunto
    \[
        A=\{x\in\mathbb R^n : \|x\|>1\}
    \]
    é aberto.
\end{exercise}

\begin{exercise}
    Considere $A=[0,+\infty[\subseteq\mathbb R$ com a métrica herdada da distância usual de $\mathbb R$.
    Determine quais dos seguintes subconjuntos são abertos no subespaço $A$.
    Justifique suas respostas.
    \begin{enumerate}[label=(\alph*)]
        \item $A$;
        \item $[0,1[$;
        \item $]0,1[$;
        \item $\{0\}$;
        \item $[1,2[$.
    \end{enumerate}
\end{exercise}

\begin{exercise}
    Considere o conjunto $\mathbb Z$ munido da métrica usual.
    Determine quais subconjuntos de $\mathbb Z$ são abertos em $\mathbb Z$.
\end{exercise}

\begin{exercise}
    Seja $M$ um conjunto munido da métrica discreta, isto é,
    \[
        d(p,q)=
        \begin{cases}
            0, & \text{se } p=q,\\
            1, & \text{se } p\neq q.
        \end{cases}
    \]
    Mostre que todo subconjunto de $M$ é aberto.
\end{exercise}